
    \chapter{Histórico de Planeamento e Riscos}
    Durante a fase de planeamento foram identificados riscos técnicos e operacionais que continuam relevantes para a evolução do projeto.

    \section{Riscos}

    Durante o desenvolvimento do projeto foram identificados vários riscos técnicos e operacionais. Estes riscos influenciaram decisões de desenho, prioridades de implementação e estratégias de validação.

    \begin{itemize}
        \item \textbf{Curva de aprendizagem tecnológica:}
        O projeto exigiu a utilização de tecnologias e conceitos com elevada complexidade, nomeadamente Kotlin Multiplatform, arquitetura hexagonal, APIs de inferência de diferentes fornecedores, autenticação baseada em OIDC/OAuth2 e o Model Context Protocol. Esta diversidade tecnológica representou um risco inicial, porque aumentou o tempo necessário para compreender os contratos externos, configurar o ambiente de desenvolvimento e definir uma arquitetura estável.
        A mitigação passou pela realização de protótipos incrementais, pelo estudo isolado das APIs de OpenAI, Anthropic e Gemini, e pela separação progressiva do sistema em módulos com responsabilidades claras.

        \item \textbf{Instabilidade de serviços externos:}
        A Gateway depende de serviços de terceiros, como APIs de inferência, servidores de autorização e ferramentas externas. Estes serviços podem apresentar indisponibilidade temporária, alterações nos contratos, limites de utilização ou comportamentos diferentes dos documentados. Este risco poderia comprometer a validação da solução e tornar os testes dependentes de fatores externos.
        Para reduzir este impacto, foram usados testes locais, servidores simulados, mocks e abstrações sobre o transporte HTTP. Além disso, a arquitetura inclui mecanismos de resiliência, como \textit{fallback} e \textit{circuit breaker}, que permitem lidar melhor com falhas de fornecedores.

        \item \textbf{Custos de utilização:}
        A execução de testes com modelos pagos pode gerar custos inesperados, especialmente em cenários de \textit{streaming}, chamadas repetidas ou testes de integração. Este risco é relevante porque a Gateway incentiva a experimentação com múltiplos fornecedores e modelos.
        A mitigação passou por limitar o número de chamadas reais, utilizar fornecedores com quotas gratuitas sempre que possível, recorrer a testes locais e validar grande parte da lógica através de contratos, adapters e servidores simulados.

        \item \textbf{Complexidade de normalização entre fornecedores:}
        As APIs de OpenAI, Anthropic e Gemini partilham conceitos semelhantes, como mensagens, modelos, parâmetros de geração e ferramentas, mas representam esses conceitos de formas diferentes. Alguns elementos, como \textit{tool calling}, \textit{streaming}, multimodalidade e metadados específicos, não têm correspondência direta entre fornecedores.
        Este risco poderia levar a uma abstração demasiado rígida ou à perda de informação específica de cada API. Para o mitigar, foi definido um domínio comum extensível, complementado por campos de opções específicas do fornecedor. Desta forma, o SDK consegue representar o núcleo comum sem impedir a passagem de características particulares de cada API.

        \item \textbf{Segurança e dados sensíveis:}
        A centralização dos pedidos numa Gateway aumenta a responsabilidade sobre autenticação, autorização, gestão de tokens, logs, métricas e chaves de API. Um erro nesta camada poderia expor dados sensíveis ou permitir acesso indevido a modelos e recursos.
        A mitigação passou pela introdução de uma pipeline de segurança com autenticação e autorização separadas, validação de tokens, integração com um Authorization Server e cuidado na recolha de métricas e logs. Ainda assim, este risco continua relevante para uma evolução futura em direção a ambientes de produção e cenários multi-tenant.

        \item \textbf{Âmbito e complexidade funcional:}
        O domínio das Gateways de IA inclui várias funcionalidades possíveis, como roteamento por custo, políticas avançadas de autorização, execução automática de ferramentas, suporte completo a MCP, cache semântica e gestão de quotas por utilizador. Implementar todas estas capacidades numa primeira versão aumentaria significativamente o risco de atrasos e instabilidade.
        A mitigação passou pela definição de um âmbito inicial mais controlado, focado na arquitetura base, nos adapters principais, na pipeline de interceptors e numa Gateway funcional construída sobre o SDK. As funcionalidades mais avançadas foram identificadas como trabalho futuro.
    \end{itemize}

    \section{Planeamento}
    \begin{table}[H]
        \centering
        \footnotesize
        \renewcommand{\arraystretch}{1.2}
        \begin{tabular}{|c|l|p{11.0cm}|}
            \hline
            \textbf{Semana} & \textbf{Período} & \textbf{Descrição} \\
            \hline
            1 & 18 fev -- 22 fev (parcial) & Estudo das APIs de inferência e protocolo MCP. \\
            2 & 23 fev -- 1 mar & Continuação da semana anterior e início da proposta. \\
            3 & 2 mar -- 8 mar & \textbf{Entrega Proposta do projeto} \\
            4 & 9 mar -- 15 mar & Configuração inicial da biblioteca, definição da arquitetura base e das interfaces. \\
            5 & 16 mar -- 22 mar & Integração com as APIs de inferência. Foco inicial em dois fornecedores estáveis. \\
            6 & 23 mar -- 29 mar & Integração do MCP no OmniAI Core. Invocação e registo dinâmico de ferramentas, gerindo o ciclo de \textit{tool calling}. \\
            7 & 30 mar -- 5 abr & Implementação de despacho configurável de modelos e mecanismos de redundância (\textit{fallback} e \textit{retry}). \\
            8 & 6 abr -- 12 abr & Implementação dos mecanismos de segurança, como controlo de acesso, preparação do contexto e remoção de dados sensíveis. \\
            9 & 13 abr -- 19 abr & Conclusão das implementações anteriores. \\
            10 & 20 abr -- 26 abr & \textbf{Apresentação de Progresso} \\
            11 & 27 abr -- 3 mai & Início do desenvolvimento da aplicação final OmniAI Gateway. Exposição da WEB API. \\
            12 & 4 mai -- 10 mai & Continuação da semana anterior e preparação das abstrações MCP para integração futura. \\
            13 & 11 mai -- 17 mai & Término de todas as APIs de comunicação para o exterior. \\
            14 & 18 mai -- 24 mai & Testes de integração do Core e da Gateway. \\
            15 & 25 mai -- 31 mai & \textbf{Entrega da Versão beta} \\
            16 & 1 jun -- 7 jun & Refinamento de testes a todos os módulos. \\
            17 & 8 jun -- 14 jun & Refinamento de código, desempenho e limpeza de bugs do sistema. \\
            18 & 15 jun -- 21 jun & Continuação da semana anterior e refinamento do relatório final. \\
            19 & 22 jun -- 28 jun 2026 & Conclusão do relatório final. \\
            20 & 29 jun -- 5 jul 2026 & Vistoria completa a todo o sistema e polimento de falhas para a entrega da Versão final. \\
            21 & 6 jul -- 11 jul 2026 (parcial) & \textbf{Entrega Final do Projeto e Relatório} \\
            \hline
        \end{tabular}
        \caption{Planeamento semanal do projeto}
    \end{table}


